% !TeX root = ../papers/nc-p3.tex
\section{階数15の全座標端での六次経路}
\label{sec:rank-fifteen-coordinate-ends}
\begin{theorem}\label{thm:rank-fifteen-coordinate-ends}
正の固定曲線上の通常な階数15窓で、$a,b,e,f$ の少なくとも一つが零ならば、
全始点の長さ6の経路次元は84で、四次中心の文字次元は $(3,2,0,0)$ である。
\end{theorem}
\begin{proof}
定理\ref{thm:rank-fifteen-word-symmetries}の反転は座標端の和集合を保つので
第一表示を調べる。定理\ref{thm:rank-fifteen-covers}の $P(n)$ と標準座標を用いる。
$a=0$ または $e=0$ は $n=\pm1$ と同値である。二位置移動で $n=1$ とすれば
$z(k^2-m^2)=2$、すなわち $(z+1)(\lambda^2z+1)=8\lambda z$ となる。
$t=z$ と置くと、窓を与える楕円曲線と座標は
\[
 Y^2=-t^3+14t^2-t,\qquad \lambda=\frac{4t+Y}{t(t+1)},\qquad
 (a,b,e,f)=\left(2,\frac{t+1}{2t},0,\frac4{t+1}\right).
\]
一方 $b=0$ または $f=0$ のとき $m\ne0$ なので $k\ne0$ で、
$n=\pm m/k$ を $P(n)=0$ に入れると $m^2=-k^2$ を得る。
逆も成立する。二位置移動と $i^2=-1$ の根の選択により $m=ik$, $n=i$ とする。
すると $f=0$、$\lambda^2z+1=-i\lambda(z+1)$ で、今度は $t=\lambda/i$ により
\[
 z=\frac{1-t}{t(t+1)},\qquad \lambda=it,\qquad
 (a,b,e,f)=(1+i,2k,1-i,0).
\]
どちらの表示にも通常点で分母の零点はない。第二の場合 $t=-1$ は倍率式に反する。

各曲線の関数体で経路基底を長さ5まで構成し、
前節の四挿入行列の階数と、四つの文字ごとに全位置で矢と可換になる四次の元の空間を調べる。\path{computations/verify_rank_fifteen_coordinate_ends.py}
は全除算で逆元を取った分子のノルムを $\mathbb Q[t]$ 上で因数分解する。
被約表示の分母の新しい零点は、初期座標の分母または既に逆にした元の零点からしか生じない。
同じ長さの疎な行を同じ順で選び、主元の列を最小と最大に変えた二回の計算で
各々挿入の階数 $36,34,36,34$ と、中心の四つの核次元 $3,2,0,0$ を得る。表の数は相異なる既約因子の個数である。
\[
\begin{array}{c|rr|r|r}
 &\text{最小列}&\text{最大列}&\text{共通}&\text{通常域の残余}\\\hline
 n=1 &56&50&17&11\\
 m=ik,\ n=i &74&56&13&7
\end{array}
\]
楕円表示では $t=0,\pm1$ と $t^2-14t+1,t^2-6t+1,t^2+10t+1$ の零点は
非通常または表示外である。二根の交換は $\lambda\mapsto1/(t\lambda)$ であり、
$\lambda^2-1$ と $(t\lambda)^2-1$ の零点を交換する。
有理表示では $t=0,\pm1$ と $t^2+1,t^2\pm2t-1$ の零点を同様に除く。
階数14と15が同じ $t$ の上に来る場合は両方を計算に含めた。

残る各既約因子に対しては101,103,109,157のいずれかの素体上の単根を一つ指定し、
$Y$ または $i$ の二つの相異なる値の双方で経路基底・四階数・中心の核次元を再計算する。
単根での良い剰余は標数零の根に持ち上がる。二つの値を調べたので、根の体への
二次拡大が分解する場合も全ての幾何点を覆う。他の根へは共役で移る。
以上で全通常域において目標階数が従う。系\ref{cor:rank-fifteen-sixth-test}と
反転・二位置移動から結論を得る。
\end{proof}
\begin{corollary}
表の最後の列の18因子上では、対象の全ての通常窓が標準Koszul左右AS正則である。
\end{corollary}
\begin{proof}
上記コードは同じ全剰余点において四次中心の文字次元 $(3,2,0,0)$ と
四中心商の全四始点の長さ13までの次元も検算する。経路・中心の次元の下界と
命題\ref{prop:fixed-residual-regularity}の中心の持上げにより、
これらの閉点では全次数の正則性が従う。表の通常域の残余以外に対する
本節の一様な結論は六次の経路と四次中心に関するものである。
\end{proof}
